% Research cards: CAR, MMG
% Checked on the web on September 23, 2026. "(to be confirmed)" marks what was not verified in a source.

\section{Load, demand and MMGD}

\subsection*{Shared sources}

\begin{table}[H]
\centering\small
\caption{Sources shared by the CAR and MMG modules.}
\begin{tabularx}{\textwidth}{P{3.2cm} L P{4.2cm}}
\toprule
\textbf{Source} & \textbf{What it provides} & \textbf{Pitfalls} \\
\midrule
BDGD (ANEEL) [1,2] & annual geodatabase by distribution utility, reference date December 31, published about 8 months after the reference date in the 2025 edition. In the unit tables: feeder, transformer, substation, municipality, class, CNAE, curve typology, monthly energy and demand, \textbf{monthly DIC and FIC}, DG code and coordinates & the open CSVs only carry \textbf{legal entities}; residential requires the full geodatabase. DG energy as injection and power in kW and W mixed: observed empirically in a subset, not confirmed in the specification \\
MMGD registry (ANEEL) [4,5] & monthly; municipality, substation and approximate coordinates, modality (local, remote, shared), class, source, power in kW & power is the sum of the \textbf{modules} (DC), not of the inverter; there is no feeder, only substation \\
CTR typical curves (ANEEL) [6] & typical consumer and typical network from the measurement campaigns of the tariff reviews, by day type, in 15-minute intervals & unit of the value not stated; each distribution utility has one campaign year; curves predate the MMGD expansion \\
ONS, verified and scheduled load [7,8] & half-hourly by load area, with components, including the \textbf{share served by MMGD} & breaks in 2021 (global load) and 2023 (estimated MMGD enters the load) [9]; ONS advises not using years before 2022 [10] \\
IBGE, 2022 Census [11,12] & aggregates and census tract mesh & 2010 and 2022 tracts are not comparable without harmonization \\
Nighttime lights (NASA Black Marble) [13,14] & 500~m, daily, monthly and annual & public lighting is not consumption; urban saturation \\
Google Open Buildings 2.5D Temporal [15] & building presence, count and height, annual from 2016 to 2023, covers Brazil & raster, no polygon; ends in 2023 \\
Senatran, fleet [16,17] & monthly by municipality and fuel since 2003 & fleet registered at the owner's domicile; rental companies concentrate registrations \\
EPE [18--21] & PDE 2035 (MMGD, electromobility, demand); the 4MD model and the R package \emph{epe4md} & national or by distribution utility; the disaggregation is the user's \\
\bottomrule
\end{tabularx}
\end{table}

% ---------------------------------------------------------------------------
\subsection{CAR \textperiodcentered{} Load and demand}

\begin{ficha}{CAR-01 \textperiodcentered{} Forecast load by substation and feeder}{V2}
\campo{Data} Substation and feeder metering at 15 minutes (client);
weather forecast and reanalysis; calendar; BDGD topology; MMGD by
feeder (MMG-01).
\campo{Approaches} Baseline: similar profile and regression with
temperature and calendar (Tao Vanilla style) [22]. Recommended: one
global quantile boosting model for all feeders (P10, P50, P90),
followed by hierarchical reconciliation transformer, feeder,
substation, area with MinT [23]; probabilistic version as in Ben Taieb et
al.\ with smart meters [24]. Recent literature: TFT [26] and foundation
models [27--30]; \textbf{Chronos-2} accepts native covariates,
needed for weather [28]; use as challenger and for feeders with
little history [31,32].
\campo{Metrics} Pinball and CRPS [33]; P10--P90 coverage (expected 80\%);
WAPE of the P50 by horizon; zero-sum residual after reconciling.
\campo{Pitfalls} Switching operations and transfers create false steps (requires
a switching log); what is metered is net load; MinT with thousands of series needs
covariance with shrinkage.
\end{ficha}

\begin{ficha}{CAR-02 \textperiodcentered{} Forecast transformer loading}{V2}
\campo{Data} BDGD: monthly energy and curve typology by unit, link
to the transformer and rated power [1]; CTR curves [6]; weather; sample
metering of the client's transformers, when available.
\campo{Approaches} Baseline: the \textbf{regulatory method} for technical
losses, which ANEEL has run in OpenDSS since 2014 from the BDGD, with
monthly energy × typical curves in 24 levels [34,35]. Recommended: the
baseline plus a calibration factor learned on metered transformers
(boosting with class composition, number of units, weather and MMGD), to
reduce the underestimation of diversity; output as peak distribution and
probability of exceeding 100\% or 120\%. Conversion to OpenDSS:
\emph{bdgd2dss} [36].
\campo{Metrics} Peak error on metered transformers; AUC-PR of
overload; Brier.
\campo{Pitfalls} CTR curves old and without MMGD [3]; wrong unit to
transformer link in the BDGD; rated power is not thermal capacity.
\end{ficha}

\begin{ficha}{CAR-03 \textperiodcentered{} Forecast peak demand}{V2}
\campo{Approaches} Baseline: historical maximum × growth.
Recommended: peak density forecast by simulation, feeding the
hourly model with many years of resampled weather and extracting the
10, 50 and 90\% exceedance probabilities (Hyndman and Fan [37]; Hong, Wilson
and Xie [38]). Recent literature: extreme values (GEV, GPD) on residuals.
\campo{Metrics} Pinball at the high quantiles; error of the peak value and
time; hit rate of ``peak day''.
\campo{Pitfalls} With MMGD the net peak shifts to the early evening;
few extremes in the history.
\end{ficha}

\begin{ficha}{CAR-04 \textperiodcentered{} Spatialized demand}{V2}
\campo{Data} Consumption by unit with coordinates (full BDGD or
client); 2022 Census tracts; Open Buildings 2.5D; Black Marble.
\campo{Approaches} Baseline: direct aggregation of units into a hexagonal
grid or tract, where there are coordinates. Recommended: dasymetric in
WorldPop style, with a random forest that predicts consumption density from built
area and volume, population, nighttime light and land use, redistributing the
known total with a sum constraint [39,40]. Recent literature:
downscaling with deep learning at about 130~m [41].
\campo{Metrics} Spatial cross-validation (municipality left out)
against the actual aggregation; exact mass conservation.
\campo{Pitfalls} Nighttime light measures lighting; \textbf{LGPD}: aggregate with
a minimum number of units per cell.
\end{ficha}

\begin{ficha}{CAR-05 \textperiodcentered{} Climate sensitivity of load}{V3}
\campo{Approaches} Baseline: piecewise regression with degree days; the
slope above the balance point (MW/°C) is the indicator. Recommended:
GAM with temperature spline, hour × temperature interaction, lagged
effects from 24 to 72~h and Bayesian hierarchy across feeders. State of the
art: distributional regression [42]. Context: the IEA estimates +3.6~GW of demand
per +1~°C in Brazil [43,44].
\campo{Pitfalls} Estimate on gross load (a hot, sunny day reduces
the net load); air-conditioning penetration changes the sensitivity over the years.
\end{ficha}

\begin{ficha}{CAR-06 \textperiodcentered{} Long-term load growth}{V3}
\campo{Approaches} Baseline: PDE rate applied top-down
[45]. Recommended: structural model by class with scenarios, reconciled with
the PDE at the top, and spatial allocation by built-area growth;
large point loads (data centers, industry) as events.
\campo{Metrics} Rolling-origin backtest over 3 to 10 years; spatial
allocation error; scenario coverage.
\end{ficha}

\begin{ficha}{CAR-07 \textperiodcentered{} Consumption profiles by area}{V4}
\campo{Approaches} Without smart metering, synthetic profile from the class
composition and CTR curves; with metering, clustering (k-medoids with DTW, GMM)
evaluated by Chicco's indices [46].
\campo{Pitfalls} A synthetic profile only reproduces the existing typologies;
LGPD in small aggregations.
\end{ficha}

\begin{ficha}{CAR-08 \textperiodcentered{} Forecast load by submarket}{cond}
\campo{Approaches} Baseline: the ONS scheduled load [8].
Recommended: boosting or TFT by subsystem with load-weighted weather and the
MMGD component modeled separately [7]; reconciliation area, subsystem, SIN.
\campo{Pitfalls} Breaks of 2021 and 2023 [9]; retroactive revisions.
\end{ficha}

% ---------------------------------------------------------------------------
\subsection{MMG \textperiodcentered{} MMGD and distributed resources}

\begin{ficha}{MMG-01 \textperiodcentered{} Installed MMGD by feeder}{V2}
\campo{Approach} (1) Join registry and BDGD by the DG code; on failure,
spatial join within the same substation. (2) Normalize units by
comparing with the registry power (a ratio close to 1,000 indicates W). (3)
For the network what counts is the plant location, not the unit that receives the credit
(remote, shared). (4) Convert modules (DC) to inverter (AC) with the
sizing factor used by EPE [20].
\campo{Metrics} Join rate; power divergence between sources by
feeder.
\campo{Pitfalls} BDGD with 8 to 20 months of lag and monthly registry: use the
registry for counts and the BDGD for topology; approximate coordinates lead
to the wrong feeder in dense areas.
\end{ficha}

\begin{ficha}{MMG-02 \textperiodcentered{} Forecast distributed generation}{V2}
\campo{Approaches} Baseline: installed power × irradiance × PR (0.75
local and 0.80 remote, 4MD assumptions) [20]. Recommended: physical model with
pvlib and upscaling by reference systems, with quantile
post-processing calibrated on the feeder metering. ONS estimates MMGD with
the ANEEL capacity, INPE irradiation and its own capacity factor [9].
\campo{Generation and injection} Generation is not injection: injection is generation minus
simultaneous self-consumption, which requires a model of the unit's load.
\campo{Pitfalls} Unknown orientation; unregistered capacity;
inverter limiting.
\end{ficha}

\begin{ficha}{MMG-03 \textperiodcentered{} Gross load × MMGD decomposition}{V2}
\campo{Approaches} Baseline: add the estimated generation to the net load.
Recommended: joint estimation, net = load(weather, calendar) $-\,k$ ×
PV proxy(irradiance), with $k$ (effective power) by regression or Kalman
filter; the estimated $k$ audits the registry and feeds MMG-05. References:
SunDance [48], unsupervised disaggregation [49], PNNL [50], game
theory [51], solar exemplars [52].
\campo{Pitfalls} Collinearity between irradiance, temperature and
air conditioning; $k$ changes over time.
\end{ficha}

\begin{ficha}{MMG-04 \textperiodcentered{} Forecast MMGD growth}{V2}
\campo{Reference model} EPE's 4MD: Bass with dynamic potential
market, 52 distribution utilities, segments residential local, other LV, HV and
remote; potential market by income; adopting fraction dependent on payback;
available as the R package \emph{epe4md} [20,21]. PDE 2035: 61.4 to 97.8~GW of
MMGD in 2035 [18].
\campo{Approaches} Baseline: \emph{epe4md} by distribution utility, allocated
by current share. Recommended: hierarchical Bass (distribution utility,
municipality, feeder) with spatial covariates; in Minas Gerais, wage,
higher education and urbanization have the largest effects and there is spatial spillover [53,54].
Recent literature: agents and system dynamics [55]; DeepSolar [56].
\campo{Pitfalls} Regulatory shocks (Lei 14.300/2022, REN 1.098/2024) that
the classic Bass does not capture, as EPE acknowledges [20].
\end{ficha}

\begin{ficha}{MMG-05 \textperiodcentered{} Unregistered MMGD}{V3}
\campo{Approaches} Baseline: persistent drop in consumption without a DG
code and inverted seasonality, with change point and correlation with
irradiance (Zhang and Grijalva 2016) [57]. Recommended: this screening plus
confirmation by computer vision on the candidates, with Brazilian
precedents [58,59]; at the feeder, the MMG-03 residual. Recent literature: Bayesian
fusion of registry and remote sensing [60]; attention to the low
reliability of detectors outside the training scale and region [61].
\campo{Pitfalls} Consumption drop from other causes; an indication is not proof;
aggregated display.
\end{ficha}

\begin{ficha}{MMG-06 \textperiodcentered{} Hosting capacity by feeder}{V3}
\campo{Context} REN 1.098/2024 created cases of exemption from the reverse
power flow analysis and an order of solutions before denying access (power
reduction, time restriction, change of connection point, works) [63,64].
\campo{Approaches} Deterministic: increase PV in the worst case until a violation.
Stochastic: Monte Carlo over location and size; Torquato et al.\ (Unicamp,
2018), on 50,000 real low-voltage networks, identified overvoltage as
the most restrictive limit and approximately lognormal hosting [65]. Time
series over 8,760~h [66]. \textbf{Recommended for scale}: OpenDSS on a
stratified sample of feeders and a surrogate that predicts hosting from
feeder attributes [67]. Recent literature: PINN [68].
\campo{Metrics} Surrogate error against OpenDSS and hit rate of the limiting
factor; agreement with access assessments.
\campo{Pitfalls} Inconsistent phases and impedances in the BDGD; unknown
primary tap; publishing ``headroom'' has regulatory and commercial effect.
\end{ficha}

\begin{ficha}{MMG-07 \textperiodcentered{} Reverse power flow and overvoltage risk}{V3}
\campo{Approaches} Baseline: P90 of injection over P10 of daytime load,
by transformer and feeder. Recommended: joint sampling of load and
PV with a copula and voltage sensitivity from the MMG-06 surrogate, calibrated
with the client's complaints and measurements.
\campo{Pitfalls} Minimum load of an unmetered transformer comes from outdated
curves; underreported events.
\end{ficha}

\begin{ficha}{MMG-08 \textperiodcentered{} Forecast electric vehicle adoption}{V4}
\campo{Data} Senatran fleet [16]; PDE 2035: electric vehicles at 23\% of
light-vehicle registrations in 2035 and demand of 0.6 to 7.8~TWh between 2025 and 2035
[19].
\campo{Approaches} Hierarchical Bass by municipality constrained to the PDE total,
separating plug-in electric vehicles from hybrids; load by vehicles × km × kWh/km ×
charging profile, allocated by CAR-04.
\campo{Pitfalls} Fleet at the owner's domicile; little measured charging
data in Brazil.
\end{ficha}

\subsection*{References for this section}
{\small
\begin{enumerate}[label={[\arabic*]},leftmargin=2.2em,itemsep=1pt]
\item ANEEL, BDGD: \url{https://dadosabertos.aneel.gov.br/dataset/base-de-dados-geografica-da-distribuidora-bdgd}
\item ANEEL, PRODIST Module 10 (REN 956/2021)
\item SBSE, impact of DG on load typologies
\item ANEEL, MMGD projects: \url{https://dadosabertos.aneel.gov.br/dataset/relacao-de-empreendimentos-de-geracao-distribuida}
\item ANEEL, DG data dictionary v2.3 (November 2025)
\item ANEEL, CTR Curva de Carga: \url{https://dadosabertos.aneel.gov.br/dataset/ctr-curva-de-carga}
\item ONS, verified load: \url{https://dados.ons.org.br/dataset/carga-energia-verificada}
\item ONS, scheduled load: \url{https://dados.ons.org.br/dataset/carga-energia-programada}
\item ONS, PMO starts to incorporate the MMGD load (2023)
\item ONS, MMGD forecasting roadmap (2024--2028)
\item IBGE, 2022 Census, aggregates by tract
\item IBGE, census tract mesh
\item NASA, Black Marble User Guide v1.2
\item NASA LAADS, VNP46A2
\item Google, Open Buildings 2.5D Temporal: \url{https://sites.research.google/gr/open-buildings/temporal/}
\item Senatran, vehicle fleet: \url{https://www.gov.br/transportes/pt-br/assuntos/transito/conteudo-Senatran/frota-de-veiculos-2025}
\item Base dos Dados, traffic statistics
\item EPE, PDE 2035, MMGD and batteries volume
\item MME and EPE, PDE 2035, electromobility volume
\item EPE, NT DEA-SEE 010/2020, 4MD model
\item EPE, R package epe4md: \url{https://epe-gov-br.github.io/epe4md/}
\item Hong et al.\ (2016), GEFCom2014, \emph{IJF} 32(3)
\item Wickramasuriya, Athanasopoulos and Hyndman (2019), MinT, \emph{JASA} 114(526)
\item Ben Taieb, Taylor and Hyndman (2021), \emph{JASA} 116(533)
\item Hong and Fan (2016), tutorial on probabilistic load forecasting, \emph{IJF}
\item Lim et al.\ (2021), TFT, \emph{IJF} 37(4)
\item Ansari et al.\ (2024), Chronos: \url{https://arxiv.org/abs/2403.07815}
\item Chronos-2 (2025): \url{https://arxiv.org/abs/2510.15821}
\item Das et al.\ (2024), TimesFM: \url{https://arxiv.org/abs/2310.10688}
\item Woo et al.\ (2024), Moirai: \url{https://arxiv.org/abs/2402.02592}
\item Emami, Sahu and Graf (2023), BuildingsBench: \url{https://arxiv.org/abs/2307.00142}
\item Foundation models on residential load (2024): \url{https://arxiv.org/abs/2410.09487}
\item Gneiting and Raftery (2007), \emph{JASA} 102(477)
\item Freitas (Unicamp), technical losses at ANEEL
\item Norven, load curves in the calculation of technical losses
\item bdgd2dss: \url{https://pypi.org/project/bdgd2dss/}
\item Hyndman and Fan (2010), \emph{IEEE TPWRS} 25(2)
\item Hong, Wilson and Xie (2014), \emph{IEEE TSG} 5(1)
\item Stevens et al.\ (2015), \emph{PLoS ONE} 10(2)
\item Consumption on a 1~km grid in 280 Chinese cities
\item Downscaling by nighttime lights, \emph{J.\ Cleaner Production} (2026)
\item Nonlinear sensitivity to temperature: \url{https://arxiv.org/abs/2503.07213}
\item IEA, Cooling a hotter world
\item Residential cooling in Brazil, \emph{Energy \& Buildings}
\item EPE, PDE 2035, electricity demand volume
\item Chicco (2012), \emph{Energy} 42(1)
\item ONS, daily energy load
\item Chen and Irwin (2017), SunDance, e-Energy
\item Unsupervised PV disaggregation, \emph{Applied Energy} (2024)
\item PNNL, net load with disaggregated PV
\item Disaggregation by game theory: \url{https://arxiv.org/abs/1907.06747}
\item Solar exemplars: \url{https://arxiv.org/abs/2009.00734}
\item DG in Minas Gerais, \emph{Energy Economics} (2025)
\item Solar adoption in Brazil, \emph{ERSS} 89 (2022)
\item System dynamics of residential solar in Brazil, \emph{RSER} (2025)
\item Yu et al.\ (2018), DeepSolar, \emph{Joule} 2(12)
\item Zhang and Grijalva (2016), \emph{IEEE TSG} 7(5)
\item PV plants in Brazil by semantic segmentation, \emph{Energies} 14(10)
\item UFMG, photovoltaic panels in orthophotos
\item Bayesian audit of PV registries: \url{https://arxiv.org/abs/2609.16294}
\item Reliability of PV detectors: \url{https://arxiv.org/abs/2408.07828}
\item SENDI, hosting capacity on a real BDGD feeder
\item eixos, reverse power flow and connection denial
\item ANEEL, simplification of microgeneration connection (2024)
\item Torquato et al.\ (2018), \emph{IEEE TPWRD} 33(2)
\item Mulenga, Bollen and Etherden (2020), \emph{IJEPES} 115
\item Qammar et al.\ (2024), \emph{IET GTD}, doi:10.1049/gtd2.12933
\item PINN for AC power flow in distribution: \url{https://arxiv.org/abs/2409.09466}
\end{enumerate}}
